\documentclass[11pt]{article}

\usepackage{amsmath, amssymb}
\usepackage{geometry}
\usepackage{listings}
\usepackage{xcolor}
\usepackage{hyperref}
\geometry{margin=1in}

\definecolor{codebg}{rgb}{0.96,0.96,0.97}
\definecolor{codeframe}{rgb}{0.8,0.8,0.82}

\lstset{
  basicstyle=\ttfamily\small,
  backgroundcolor=\color{codebg},
  frame=single,
  rulecolor=\color{codeframe},
  breaklines=true,
  columns=fullflexible,
  showstringspaces=false,
  xleftmargin=1em,
  xrightmargin=1em,
  aboveskip=1em,
  belowskip=1em,
}

\title{Blueboard Phase 2, explained}
\author{}
\date{}

\begin{document}
\maketitle

\section*{What this document is}
Phase 2 (\texttt{PLAN.md}, tag-in-progress \texttt{v0.2-vector-pen}) is where Blueboard
went from "a camera over an empty grid" to an actual whiteboard: you can draw, and what
you draw is a real vector object, not a bitmap. It took more iteration than Phase 1 —
three separate rounds of bug-hunting, each found by ruba comparing blueboard's rendering
directly against Miro's, plus a round of crash-resilience work. This note walks through
all of it in the order it actually happened, including the Rust syntax and the research
that shaped each fix. Like the Phase 1 note, it's a companion to the code, not a
replacement — file references point at the real source.

\tableofcontents

\section{The shape of the problem}
Drawing a pen stroke well requires four things to happen in order:
\begin{enumerate}
  \item \textbf{Catch pointer input} without recording more noise than signal.
  \item \textbf{Fit a smooth curve} through those points — a real vector object, per the
    brief, not a pixel trace.
  \item \textbf{Store it} in a way that supports fast "what's on screen right now"
    queries, since boards are meant to scale.
  \item \textbf{Render it} as anti-aliased GPU triangles that look as smooth as what you
    see in a drawing app you'd actually trust.
\end{enumerate}
Each step turned out to have a non-obvious failure mode, found only by looking at the
rendered result next to Miro's — not by reading the code. That pattern (write it, compare
against a real reference, fix what the comparison reveals) is the throughline of this
whole phase.

\section{\texttt{blueboard-core}: fitting and the scene graph}

\subsection{Reusing a curve fitter instead of writing one}
File: \texttt{crates/blueboard-core/src/fitting.rs}. \texttt{planning/math.tex} derives
Schneider's classic curve-fitting algorithm (chord-length parameterize, least-squares fit,
subdivide on error) as the theoretical basis for turning pointer samples into a
\texttt{BezPath}. Implementing that by hand was the original plan — but \texttt{kurbo}
0.13 turned out to ship its own general curve-fitting module
(\texttt{kurbo::fit\_to\_bezpath}), built around a trait:
\begin{lstlisting}
pub trait ParamCurveFit {
    fn sample_pt_tangent(&self, t: f64, sign: f64) -> CurveFitSample;
    fn sample_pt_deriv(&self, t: f64) -> (Point, Vec2);
    fn break_cusp(&self, range: Range<f64>) -> Option<f64>;
}
\end{lstlisting}
Rather than reimplement Schneider's method, \texttt{fitting.rs} implements this trait for
a small \texttt{Polyline} type — the raw pointer samples, treated as a piecewise-linear
curve parameterized by normalized cumulative arc length — and hands it to kurbo's own
fitter. Same underlying math as \texttt{math.tex} (arc-length parameterization,
error-bounded subdivision), but using a maintained implementation instead of a bespoke
one, per the project's own ``reuse existing utilities'' convention.
\begin{lstlisting}
fn sample(&self, s: f64, sign: f64) -> (Point, Vec2) {
    let mut i = match self
        .cumulative
        .binary_search_by(|len| len.partial_cmp(&s).unwrap())
    {
        Ok(i) => i,
        Err(i) => i.saturating_sub(1),
    };
    i = i.min(self.points.len() - 2);
    if sign < 0.0 && i > 0 && (self.cumulative[i] - s).abs() < f64::EPSILON {
        i -= 1;
    }
    let (p0, p1) = (self.points[i], self.points[i + 1]);
    /* ... */
}
\end{lstlisting}
\texttt{binary\_search\_by} is the standard library's binary search, generalized with a
comparator closure since \texttt{f64} doesn't implement \texttt{Ord} (NaN breaks total
ordering, so Rust makes you say explicitly how to compare). It returns
\texttt{Ok(index)} on an exact match or \texttt{Err(insertion\_point)} otherwise — the
\texttt{match} handles both, and \texttt{saturating\_sub(1)} avoids underflow at index 0
(an unsigned integer can't go negative; this clamps instead of panicking). The \texttt{sign}
parameter exists because \texttt{ParamCurveFit} may sample exactly at a breakpoint between
two input segments, which is ambiguous — \texttt{sign} picks which side's tangent to use,
so a corner gets a well-defined direction from both sides instead of an arbitrary one.

\subsection{The scene graph and its spatial index}
File: \texttt{crates/blueboard-core/src/scene.rs}.
\begin{lstlisting}
pub struct Scene {
    next_id: u64,
    strokes: HashMap<u64, Stroke>,
    index: RTree<IndexedBounds>,
}

struct IndexedBounds {
    id: u64,
    envelope: AABB<[f64; 2]>,
}

impl RTreeObject for IndexedBounds {
    type Envelope = AABB<[f64; 2]>;
    fn envelope(&self) -> Self::Envelope {
        self.envelope
    }
}
\end{lstlisting}
The R-tree (\texttt{rstar}) doesn't index strokes directly — it indexes a tiny
\texttt{IndexedBounds} struct holding just an id and a bounding box. The actual
\texttt{Stroke} (with its full \texttt{BezPath}) lives in a separate \texttt{HashMap},
looked up by id after a spatial query. This split matters: the R-tree's internal nodes
store copies of whatever you insert, so indexing full strokes would duplicate potentially
large path data all through the tree; indexing a bounding box is cheap to copy no matter
how complex the stroke is.
\begin{lstlisting}
pub fn strokes_in_viewport(&self, min: DVec2, max: DVec2) -> impl Iterator<Item = &Stroke> {
    let envelope = AABB::from_corners([min.x, min.y], [max.x, max.y]);
    self.index
        .locate_in_envelope_intersecting(envelope)
        .filter_map(move |entry| self.strokes.get(&entry.id))
}
\end{lstlisting}
\texttt{-> impl Iterator<Item = \&Stroke>} is an \emph{existential return type}: the
caller gets ``some concrete iterator type that yields \texttt{\&Stroke}'' without that
type being named — it's actually a chain of \texttt{rstar} and \texttt{filter\_map}
adapter types nobody wants to spell out. \texttt{move} on the closure captures
\texttt{self.strokes} by reference into the closure so it outlives the function call,
which a non-\texttt{move} closure wouldn't guarantee under the borrow checker here. This
is the viewport-culling query \texttt{blueboard-render} runs every frame (\S\ref{sec:cull}).

\section{\texttt{blueboard-render}: tessellation, and bug hunt \#1}

\subsection{From \texttt{kurbo} paths to GPU triangles}
File: \texttt{crates/blueboard-render/src/lib.rs}. Two libraries, two path types: the
scene graph stores \texttt{kurbo::BezPath}, but \texttt{lyon} (the tessellator) has its
own \texttt{lyon::path::Path}. \texttt{kurbo\_path\_to\_lyon} walks one and replays it on
the other's builder:
\begin{lstlisting}
fn kurbo_path_to_lyon(path: &kurbo::BezPath) -> LyonPath {
    let mut builder = LyonPath::builder();
    let mut closed = false;
    for el in path.elements() {
        match *el {
            kurbo::PathEl::MoveTo(p) => { builder.begin(point(p.x as f32, p.y as f32)); }
            kurbo::PathEl::CurveTo(c1, c2, p) => {
                builder.cubic_bezier_to(/* ... */);
            }
            /* ... */
        }
    }
    if !closed { builder.end(false); }
    builder.build()
}
\end{lstlisting}
\texttt{path.elements()} returns an iterator over \texttt{kurbo::PathEl} — an enum with
one variant per SVG-style path command. \texttt{match *el} dereferences each element (the
iterator yields references; the variants are small \texttt{Copy} data, so dereferencing to
match by value is cheap and avoids a wall of \texttt{\&PathEl::Foo(ref p)} patterns).
Tracking \texttt{closed} matters because lyon's builder requires exactly one
\texttt{end()} call per subpath — calling it twice (once from a \texttt{ClosePath}
element, once unconditionally afterward) would panic.

\subsection{Tessellating a stroke, and letting bad strokes fail softly}
\begin{lstlisting}
fn tessellate_stroke(path: &kurbo::BezPath, width: f32, tolerance: f32, out: &mut Vec<Vertex>) {
    let lyon_path = kurbo_path_to_lyon(path);
    let options = StrokeOptions::default()
        .with_line_width(width)
        .with_tolerance(tolerance)
        .with_line_join(lyon::tessellation::LineJoin::Round)
        .with_start_cap(lyon::tessellation::LineCap::Round)
        .with_end_cap(lyon::tessellation::LineCap::Round);

    let mut buffers: VertexBuffers<Vertex, u32> = VertexBuffers::new();
    let result = tessellator.tessellate_path(&lyon_path, &options, &mut BuffersBuilder::new(
        &mut buffers,
        |vertex: StrokeVertex| Vertex { position: vertex.position().to_array() },
    ));
    if let Err(err) = result {
        log::error!("stroke tessellation failed, skipping this stroke: {err:?}");
        return;
    }
    out.extend(buffers.indices.iter().map(|&i| buffers.vertices[i as usize]));
}
\end{lstlisting}
\texttt{BuffersBuilder::new(\&mut buffers, closure)} is \texttt{lyon}'s pattern for
``build me vertex/index buffers, and here's how to turn the tessellator's internal vertex
representation into mine'' — the closure itself satisfies lyon's
\texttt{StrokeVertexConstructor} trait via a blanket impl (\texttt{lyon} implements that
trait for any \texttt{Fn(StrokeVertex) -> V}), so no explicit trait impl block is needed,
just a closure. \texttt{round} joins/caps replaced lyon's default \texttt{miter} joins —
found necessary in bug hunt \#1 below. The \texttt{if let Err(...) \{ log::error!(...);
return; \}} pattern — rather than the \texttt{.expect(...)} this originally used — is
Phase 2's crash-resilience pass (\S\ref{sec:crashes}): one malformed stroke now gets
dropped and logged instead of taking down the whole running app.

\subsection{Viewport culling}
\label{sec:cull}
\begin{lstlisting}
let board_min = camera.screen_to_board(DVec2::ZERO, viewport_f64);
let board_max = camera.screen_to_board(viewport_f64, viewport_f64);
let mut stroke_vertices = Vec::new();
for stroke in scene.strokes_in_viewport(board_min, board_max) {
    tessellate_stroke(&stroke.path, STROKE_WIDTH, tessellation_tolerance, &mut stroke_vertices);
}
\end{lstlisting}
Every frame, the camera's current viewport corners are converted to board space and used
to query the scene's R-tree (\S3.2) — only strokes actually on screen get tessellated.
\texttt{PLAN.md} calls this out explicitly as ``wired in now even though boards are
small'': the payoff doesn't show up until boards have thousands of strokes, but retrofitting
spatial culling onto a renderer that assumed ``just loop over everything'' is a much bigger
change than building it in from the start.

\subsection{Bug hunt \#1: ``the lines look pixelated''}
\label{sec:bughunt1}
ruba's first side-by-side comparison against Miro showed visibly jagged, faceted stroke
edges. Three separate causes, found by elimination:

\paragraph{Cause 1: no anti-aliasing.} The pipeline had \texttt{multisample:
wgpu::MultisampleState::default()} — sample count 1, meaning every triangle edge is a hard
binary in/out test at rasterization. Fixed with 4$\times$ MSAA: an offscreen multisampled
color texture that every pipeline renders into, resolved down into the swapchain's
single-sample texture at the end of the pass:
\begin{lstlisting}
fn create_msaa_view(device: &wgpu::Device, config: &wgpu::SurfaceConfiguration) -> wgpu::TextureView {
    let texture = device.create_texture(&wgpu::TextureDescriptor {
        size: wgpu::Extent3d { width: config.width, height: config.height, depth_or_array_layers: 1 },
        sample_count: SAMPLE_COUNT, /* = 4 */
        format: config.format,
        usage: wgpu::TextureUsages::RENDER_ATTACHMENT,
        /* ... */
    });
    texture.create_view(&wgpu::TextureViewDescriptor::default())
}
\end{lstlisting}
and in the render pass:
\begin{lstlisting}
color_attachments: &[Some(wgpu::RenderPassColorAttachment {
    view: &self.msaa_view,
    resolve_target: Some(&view),   // the real swapchain texture
    ops: wgpu::Operations {
        load: wgpu::LoadOp::Clear(/* board blue */),
        store: wgpu::StoreOp::Discard,  // only the resolved texture needs to survive
    },
    /* ... */
})],
\end{lstlisting}
GPU-side, MSAA renders each pixel at 4 sample locations instead of 1, and a triangle edge
that only covers some of those samples gets a blended (partially transparent) color at
that pixel instead of a hard on/off choice — this is literally what removes the
staircasing. \texttt{StoreOp::Discard} on the multisampled attachment is a minor
efficiency detail: once it's resolved into the single-sample swapchain texture, the
4$\times$-larger multisampled data doesn't need to be written back anywhere.

\paragraph{Cause 2: tolerance not scaled by zoom.} \texttt{planning/math.tex} \S3 is
explicit that flattening tolerance must be in screen-space pixels, dividing by zoom before
use, so detail automatically tightens when zoomed in. The curve \emph{fit} tolerance
(\S\ref{sec:bughunt2}) got this right from the start; the \emph{tessellation} tolerance
passed to lyon was a hardcoded board-space \texttt{0.25}, violating the project's own
documented design — fine at 100\% zoom (0.25 screen px), visibly coarse once zoomed in
(0.25 $\times$ zoom screen px of error). Fixed by computing it the same way the fit
tolerance already was:
\begin{lstlisting}
let tessellation_tolerance = STROKE_TESSELLATION_TOLERANCE_PX / camera.zoom as f32;
\end{lstlisting}

\paragraph{Cause 3: miter joins.} lyon's default \texttt{LineJoin::Miter} produces sharp
corners at every curve-fit segment boundary; Miro's pen reads as rounded. Switched to
\texttt{LineJoin::Round} with round caps (\S4.2) — a one-line change, visibly softer
strokes.

\subsection{Crash resilience: making wgpu's errors recoverable}
\label{sec:crashes}
Investigating ``the app sometimes crashes for no reason'' turned up a specific,
previously-seen clue: Phase 1's shader bug had logged
\texttt{"Handling wgpu errors as fatal by default"} right before panicking. That's wgpu's
actual default: any \emph{uncaptured} device error (a validation error not wrapped in an
explicit error scope) panics the whole process. wgpu provides an escape hatch:
\begin{lstlisting}
device.on_uncaptured_error(Arc::new(|error| {
    log::error!("wgpu uncaptured error: {error}");
}));
device.set_device_lost_callback(|reason, message| {
    log::error!("wgpu device lost: {reason:?}: {message}");
});
\end{lstlisting}
\texttt{Arc<dyn UncapturedErrorHandler>} — \texttt{UncapturedErrorHandler} is defined as a
blanket trait over any \texttt{Fn(Error) + Send + Sync + 'static}, so an ordinary closure
satisfies it automatically; \texttt{Arc} (not \texttt{Box}) because wgpu may need to share
the handler across internal threads. Installing these turns ``the window just vanished''
into a logged, diagnosable line — the single highest-value change in this whole bug-hunt
round, since it addresses a whole \emph{class} of potential crash, not one specific bug.

\section{\texttt{blueboard-app}: input handling, and bug hunts \#2 and \#3}

\subsection{Re-mapping input for the pen}
File: \texttt{crates/blueboard-app/src/main.rs}. Phase 1's left-drag-to-pan had to move:
left-drag now draws, middle-drag pans, scroll still zooms — a small but deliberate choice,
since the pen is the thing actually being built this phase.

\subsection{Bug hunt \#2: real mouse jitter becomes fake curvature}
\label{sec:bughunt2}
A second comparison against Miro, this time drawing a near-horizontal line specifically,
showed visible staircase kinks that \S\ref{sec:bughunt1}'s anti-aliasing fix didn't touch
— a different bug, upstream of rendering entirely. Real pointer input has a pixel or two
of jitter even moving in a straight line; every \texttt{CursorMoved} event was being
recorded and fit with a tight 1.5px tolerance, so the curve fitter faithfully reproduced
the jitter as real direction changes. Worst-case visible on a horizontal line, where the
jitter is the \emph{entire} vertical variation of the stroke.

Two changes: loosen the fit tolerance (1.5px $\to$ 2.5px, enough to smooth over a pixel or
two of noise), and \emph{decimate} raw samples by screen-space distance before they reach
the fitter:
\begin{lstlisting}
} else if self.drawing
    && cursor.distance(self.last_sample_screen) >= PEN_MIN_SAMPLE_DISTANCE_PX
{
    /* record point */
    self.last_sample_screen = cursor;
}
\end{lstlisting}
Pointer-move events arrive far more densely than meaningful hand motion; most of that
density is jitter. Dropping points closer than 2 screen-px to the last \emph{recorded}
one (not the last \emph{event}) stops near-stationary noise from ever reaching the fitter.

\subsection{Bug hunt \#3: decimation thins density, not noise}
A third comparison, zoomed in far past where the stroke was drawn, still showed small
regular kinks — decimation and a looser tolerance had helped but not solved it. The
reason: decimation only controls how many points survive, not where they are. Each
surviving point is still exactly wherever the noisy cursor was; that's invisible within
the fit tolerance at the zoom level drawn, but zooming in further effectively magnifies
the same noise.

Researching how other open-source whiteboard/drawing tools solve this led to
\href{https://github.com/steveruizok/perfect-freehand}{perfect-freehand} (Steve Ruiz),
used by both \href{https://tldraw.dev/sdk-features/draw-shape}{tldraw} and
\href{https://github.com/excalidraw/excalidraw/pull/3512}{Excalidraw} for their pen tools.
Its relevant idea is the \emph{streamline} parameter: pull each incoming point toward a
running average of recent positions — a live low-pass filter on the input stream itself,
not just its density. (perfect-freehand's full approach also builds a variable-width
outline polygon rather than a centerline-plus-width stroke, a bigger architectural
difference not adopted here — just the streamline idea.)
\begin{lstlisting}
self.smoothed_cursor = self.smoothed_cursor.lerp(cursor, PEN_STREAMLINE_ALPHA); // = 0.5
if self.smoothed_cursor.distance(self.last_sample_screen) >= PEN_MIN_SAMPLE_DISTANCE_PX {
    /* record self.smoothed_cursor, not the raw cursor */
}
\end{lstlisting}
\texttt{DVec2::lerp(self, other, t)} computes $(1-t)\cdot\text{self} + t\cdot\text{other}$
— an exponential moving average falls out of calling this every frame with the same
\texttt{t}: each update pulls the smoothed point halfway toward the new raw sample,
so old noise decays geometrically rather than being remembered forever.
\texttt{smoothed\_cursor} is snapped to the raw cursor position on pen-down (no lag at the
very start of a stroke) and is what gets decimated and fed to the fitter — the raw
\texttt{cursor} is no longer recorded directly at all. This is the fix that actually
closed the gap: re-tested the same jittery horizontal line, zoomed in to the same level
that showed kinks before, and it was smooth.

\subsection{Zoom range}
\texttt{MAX\_ZOOM} raised from the brief's original 200\% to 400\% in
\texttt{blueboard-core/src/camera.rs}, after ruba compared directly against how far Miro
lets you zoom. \texttt{MIN\_ZOOM} (1\%) unchanged.

\subsection{Logging: stderr isn't enough for a GUI app}
File: \texttt{crates/blueboard-app/src/logging.rs}. The crash-resilience work in
\S\ref{sec:crashes} converts wgpu errors into log lines — but a log line only helps if
someone can actually read it, and a GUI app's stderr is easy to lose (closed terminal,
launched from a desktop icon). Rather than reach for \texttt{env\_logger} (which only
writes to one target), a small custom \texttt{log::Log} implementation writes every line
to \emph{both} stderr and a persistent file:
\begin{lstlisting}
impl Log for FileAndStderrLogger {
    fn log(&self, record: &Record) {
        let line = format!("[{}] {:5} {}: {}", unix_timestamp(), record.level(), record.target(), record.args());
        eprintln!("{line}");
        if let Ok(mut file) = self.file.lock() {
            let _ = writeln!(file, "{line}");
        }
    }
    /* ... */
}
\end{lstlisting}
\texttt{Mutex<File>} guards the file handle since \texttt{Log} requires \texttt{Sync} (the
global logger can be called from any thread); \texttt{.lock()} returning an \texttt{Err}
(a poisoned mutex, from an earlier panic while holding the lock) is handled by just
skipping that log line via \texttt{if let Ok(...)} rather than panicking again — a logger
that can itself cause a crash would defeat the point. The log file lives at
\texttt{\$XDG\_STATE\_HOME/blueboard/blueboard.log} (falling back to
\texttt{~/.local/state/...}), matching the XDG convention already used for the config file
path decided in Phase 0's research — config and runtime state are conventionally separate
directories, not reused for both.
\begin{lstlisting}
fn install_panic_hook(path: PathBuf) {
    let default_hook = std::panic::take_hook();
    std::panic::set_hook(Box::new(move |info| {
        if let Ok(mut file) = open_log_file(&path) {
            let _ = writeln!(file, "[{}] PANIC: {info}", unix_timestamp());
        }
        default_hook(info);
    }));
}
\end{lstlisting}
\texttt{std::panic::take\_hook()} retrieves Rust's current panic hook (by default, the one
that prints a backtrace-ish message to stderr) and \emph{removes} it, handing back
ownership; the new hook wraps it, so the original behavior (stderr output) still happens
via \texttt{default\_hook(info)} at the end — this appends to the log file \emph{in
addition to}, not instead of, the normal panic message. Combined with
\S\ref{sec:crashes}'s wgpu error handlers, the two together cover both failure modes this
phase has actually seen: an uncaptured wgpu error (now logged, not fatal) and an
unexpected panic anywhere else in the program (now logged to a file that survives a closed
terminal, even though the panic itself is still fatal — turning an unexplained
disappearance into a diagnosable one).

\section{Tools and research consulted this phase}
\begin{description}
  \item[\texttt{kurbo}'s \texttt{fit} module] (\S2.1) — avoided a hand-rolled Schneider
    implementation.
  \item[\texttt{lyon\_tessellation}] (\S4.1--4.2) — stroke tessellation, read directly
    from its source (as with wgpu in Phase 1) to get the exact 1.0.19 API shape right.
  \item[\href{https://github.com/steveruizok/perfect-freehand}{perfect-freehand}]
    (\S\ref{sec:bughunt2}) — the streamline-smoothing concept, the open-source reference
    point that resolved bug hunt \#3.
  \item[wgpu's error-scope docs] (\S\ref{sec:crashes}) — \texttt{on\_uncaptured\_error} /
    \texttt{set\_device\_lost\_callback} as the documented alternative to wgpu's default
    panic-on-error behavior.
\end{description}

\section{A verification-methodology lesson}
Worth recording since it shaped how to read ``verified by running the app'' claims in this
phase's git history: background test processes were launched with
\texttt{DISPLAY=:0}, forcing the X11/XWayland backend so \texttt{xdotool} could drive
them — but \texttt{:0} is ruba's actual live desktop session, not an isolated test
display. A test run showed a stroke logged (\texttt{"stroke 0 committed from 151 raw
samples"}) in a session where no synthetic input had been sent yet — real mouse activity
on the shared display had landed on the test window. Earlier in the phase, screenshots
that appeared to show stale content from a previous session most likely had the same
explanation, not a compositor caching bug as first suspected. Background GUI testing on a
shared display can't be treated as fully isolated; a genuinely clean test needs either a
dedicated virtual display or a window confirmed to have received no outside input during
the test window.

\section{What's next}
Phase 3, per \texttt{PLAN.md}: the \texttt{.bbd} file format — serializing the scene graph
with \texttt{serde}/\texttt{bincode}, compressing with \texttt{zstd}, and round-trip tests
from the start, since (per the project's own safety conventions) file-format bugs are the
kind that can destroy a user's work.

\end{document}
